% vim: tw=78 ai spell spelllang=en_gb sw=2
% syntax spell toplevel

\documentclass[10pt]{llncs}
\usepackage{color}
\usepackage{colortbl}
\usepackage{url}
\usepackage[T1]{fontenc}
\usepackage{graphicx}
\usepackage{listings}
\usepackage{hyperref}
\usepackage{multirow}
\usepackage{tikz}
\usepackage{xspace}
\pagestyle{plain}

\usetikzlibrary{arrows, positioning, decorations.pathreplacing}

\lstset{language=C,tabsize=2,showstringspaces=false,
	breaklines=true,breakatwhitespace=true,escapeinside={(*@}{@*)},
	captionpos=b,basicstyle=\sffamily\small}
%numbers=left,frame=l

\newcommand{\code}[2][]{\lstinline[columns=fullflexible, #1]$#2$}
\newcommand{\codeasm}[2][]{\code[language={[x86masm]Assembler},#1]{#2}}

\newcommand{\xes}{\texttt{x86}\xspace}
\newcommand{\xesl}{\texttt{x86\_32}\xspace}
\newcommand{\xesh}{\texttt{x86\_64}\xspace}

\newcommand{\gcc}{\textsc{Gcc}\xspace}
\newcommand{\orc}{\textsc{Orc}\xspace}

\setcounter{tocdepth}{1}

\begin{document}

\title{Unwinding the State of Unwinder}
\author{Jiri~Slaby}
\institute{Suse Labs \\
\email{jslaby@suse.cz}}

\maketitle

\begin{abstract}
  We have been using our own implementation of DWARF unwinder in every SUSE
  release in the past decade. First, it makes the kernel faster, second, it
  prunes the backtraces from false positives. However, this is obviously an
  out-of-tree solution. Given the progress in the live patching area in
  upstream, there is a call for some reliable unwinder in upstream too. This
  talk describes the path of our unwinder to upstream, its rejection, proposed
  solution and final SLE15 implementation.
\end{abstract}

\section{Introduction} \label{sec:intro}
Most of us already met a kernel crash, \textit{BUG}, or at least a kernel
\textit{WARNING}. If the report produced by the kernel did not contain
information what functions were visited up to the point of failure, such
report would be helpless for developers. This piece of information is called a
\textit{stack trace} and an example is depicted in Fig.\
\ref{fig:stack_trace}.

\begin{figure}[htb]
  \begin{center}
%RIP: 0010:crash_now+0xae/0x170 [crasher]
%RSP: 0018:ffffb82980d1bc80 EFLAGS: 00010282
%RAX: 0000000000000023 RBX: 0000000000000000 RCX: 0000000000000000
%RDX: ffff902dd3296888 RSI: ffff902dd328e568 RDI: ffff902dd328e568
%RBP: ffffb82980d1bc88 R08: 0000000000000000 R09: 0000000000000000
%R10: 0000000000000000 R11: 0000000000000000 R12: ffffffffc04a8000
%R13: ffff902dcb5c8e20 R14: 0000000000000001 R15: ffff902d7580f880
%FS:  00007fa4fa242040(0000) GS:ffff902dd3280000(0000) knlGS:0000000000000000
%CS:  0010 DS: 0000 ES: 0000 CR0: 0000000080050033
%CR2: 0000000000000000 CR3: 0000000089138000 CR4: 00000000000006e0
%Call Trace:
    \begin{tabular}{c}
      \begin{lstlisting}[language=XML,columns=fullflexible]
crasher_init+0x4d/0x1000 [crasher](*@\tikz[remember picture,baseline=-.5ex]{\node [coordinate] (LSTSTMODINIT) {};}@*)
(*@\tikz[remember picture,baseline=-1.5ex]{\node [coordinate] (LSTSTINITTOP) {};}@*)do_one_initcall+0x50/0x1a0
do_init_module+0x5f/0x1f0
(*@\tikz[remember picture,baseline=0ex]{\node [coordinate] (LSTSTINITBOT) {};}@*)load_module+0x1694/0x1ea0
(*@\tikz[remember picture,baseline=-1.5ex]{\node [coordinate] (LSTSTMODPTOP) {};}@*)SYSC_finit_module+0xd7/0xf0
(*@\tikz[remember picture,baseline=0ex]{\node [coordinate] (LSTSTMODPBOT) {};}@*)SyS_finit_module+0xe/0x10
entry_SYSCALL_64_fastpath+0x23/0xc2(*@\tikz[remember picture,baseline=-.5ex]{\node [coordinate] (LSTSTSYSCALL) {};}@*)
      \end{lstlisting}

  \tikzstyle{ptr} = [->,>=stealth',shorten > =3pt]
  \begin{tikzpicture}[remember picture,overlay,node distance=1.5em]
    \node [right=of LSTSTMODINIT] { module's init function }
      edge[ptr] (LSTSTMODINIT);

    \node [coordinate,right=4.4cm of LSTSTINITTOP] (ITOPSHIFTED) {};
    \node [coordinate,right=4.4cm of LSTSTINITBOT] (IBOTSHIFTED) {};
    \draw[-,thick,decorate,decoration={brace,amplitude=4pt}] (ITOPSHIFTED) -- (IBOTSHIFTED)
      node[midway] (MID) {};
    \node [right=1mm of MID] { load and init module };

    \node [coordinate,right=4.7cm of LSTSTMODPTOP] (MTOPSHIFTED) {};
    \node [coordinate,right=4.7cm of LSTSTMODPBOT] (MBOTSHIFTED) {};
    \draw[-,thick,decorate,decoration={brace,amplitude=3pt}] (MTOPSHIFTED) -- (MBOTSHIFTED)
      node[midway] (MID) {};
    \node [right=1mm of MID] { \code{modprobe} system call };

    \node [right=of LSTSTSYSCALL] { system call }
      edge[ptr] (LSTSTSYSCALL);
  \end{tikzpicture}
    \end{tabular}
  \end{center}

  \caption{Example of a stack trace. Read from bottom to top.}
  \label{fig:stack_trace}
\end{figure}

A developer can take reports like these and start debugging using
disassembler, \textsc{gdb}, debug-enabled kernels and so on. As we can see,
offsets in the functions are reported too, so we can actually track precisely
which \codeasm{call}s were invoked in the chain in every function.

\subsection{Frames and (Optional) Frame Pointers} \label{sec:fp}
There are many approaches on where and how to extract a stack trace entries.
In the kernel, we use the fact that every \codeasm{call} to a function
implicitly stores a \emph{return address} to the stack before yielding the
execution to the function and every \code{return} reads and implicitly removes
this address.  So every call is said to create a \emph{stack frame} for a
particular function call. And every \code{return} destroys it.  There is an
example of code and stack in Fig.\ \ref{fig:stack_frame}. On the left side,
there are three functions which cause the program to crash by \code{abort}.
When that happens, the possible state of the stack is depicted on the right
side.

\begin{figure}[htb]
  \begin{center}
      \tikzstyle{ptr} = [->,>=stealth',shorten > =3pt,semithick]
      \tikzstyle{border} = [thick]
      \tikzstyle{leftbrace} = [-,semithick,decorate,
	decoration={brace,mirror,amplitude=5pt,raise=2pt}]
      \tikzstyle{lefttext} = [left=5pt of MID,anchor=south,rotate=90,
	align=center]
      \newlength{\hw}
      \setlength{\hw}{1.7cm}
      \begin{tikzpicture}[node distance=2mm]
	\node (PRE) { \vdots };

	\node[below=of PRE,align=center] (CFA1) { return address \\
	  from \code{fun1} (to \code{main}) };
	\node[below=of CFA1] (SRBP1) { saved \code{BP} of \code{main}};
	\node[below=of SRBP1] (X) { variable \code{xx} };

	\node[below=of X,align=center] (CFA2) { return address \\
	  from \code{fun2} (to \code{fun1}) };
	\node[below=of CFA2] (SRBP2) { saved \code{BP} of \code{fun1}};
	\node[below=of SRBP2] (Y) { variable \code{yy} };

	\node[below=of Y,align=center] (POSTCOM) { stack grows here \\
	  up to \code{abort}};
	\node[coordinate,below=6mm of POSTCOM] (POST) {}
	  edge[ptr,<-,shorten > =0pt,shorten < = 2pt] (POSTCOM);

	\node [left=\hw of PRE.north] (TL) {};
	\node [right=\hw of PRE.north] (TR) {};
	\node [below left=1mm and \hw of POST.south] (BL) {};
	\draw[border] (TL) -- (BL) node[midway] (LCEN) {};
	\node [below right=1mm and \hw of POST.south] {}
	  edge[border] (TR);

	\foreach \from/\to/\left/\right/\line in
	    {PRE/CFA1/MAINFUN1/CALLFUN1/thick, CFA1/SRBP1/D1/D2/dashed,
	    SRBP1/X/D1/ORBP/dashed, X/CFA2/FUN1FUN2/CALLFUN2/thick,
	    CFA2/SRBP2/D1/D2/dashed, SRBP2/Y/D1/RBP/dashed,
	    Y/POSTCOM/D1/RSP/thick} {
	  \path (\from) -- (\to) node[midway] (MID) {};

	  \node [coordinate,left=\hw of MID] (\left) {};
	  \node [coordinate,right=\hw of MID] (\right) {}
	    edge[semithick,\line] (\left);
	}

	\node[right=7mm of RBP] { \code{BP} }
	  edge[ptr] (RBP);
	\node[right=7mm of RSP] { \code{SP} }
	  edge[ptr] (RSP);
	\node[above right=-1mm and 8mm of CALLFUN1,align=center] { called \\
	  \code{fun1} }
	  edge[ptr,dashed] (CALLFUN1);
	\node[above right=-1mm and 8mm of CALLFUN2,align=center] { called \\
	  \code{fun2} }
	  edge[ptr,dashed] (CALLFUN2);

	\draw (SRBP2.east) edge[ptr,in=0,out=30] (ORBP);

	\node[coordinate,above=23mm of ORBP] (TMP) {};
	\draw (SRBP1.east) edge[ptr,in=0,out=30] (TMP);

	\draw[leftbrace] (TL) -- (MAINFUN1) node[midway] (MID) {};
	\node[lefttext] { \code{main}'s \\ frame };

	\draw[leftbrace] (MAINFUN1) -- (FUN1FUN2) node[midway] (MID) {};
	\node[lefttext] { \code{fun1}'s \\ frame };

	\draw[leftbrace] (FUN1FUN2) -- (BL) node[midway] (MID) {};
	\node[lefttext] { \code{fun2}'s \\ frame };

	\node[left=2.6cm of LCEN] (CODE) {
	  \begin{lstlisting}
int main()
{
  fun1();
}

void fun1(void)
{
  long xx;

  fun2();
}

void fun2(void)
{
  long yy;

  abort();(*@ \normalfont $\leftarrow$ we are here@*)
}
	  \end{lstlisting}
	};

      \end{tikzpicture}
  \end{center}

  \caption{Example of code and its possible stack. The stack grows down, so
    the events happened from the top to bottom (as well as the code).
    \code{SP} = current stack pointer, \code{BP} = current base/frame
    pointer.}
  \label{fig:stack_frame}
\end{figure}

More precisely, when \code{fun1} is called from \code{main}, the address of
the instruction right after the \codeasm{call} is stored on the stack (to know
exactly where to return) and \code{fun1} now continues the execution.
Traditionally, the first thing a function does is storing the caller's
base/frame pointer (\code{BP}) to the stack (for later restore) and making
\code{BP} to point to \code{SP} -- to the current ``bottom'' of the stack,
i.e.\ the current frame.  This happens in a so-called \emph{prologue}.  Further,
given \code{fun1} also contains a local variable \code{xx}, which is to be
allocated on stack, this is done too. Then there is a call to \code{fun2} and
the steps are repeated until \code{abort}'s body causes the program to exit.

If \code{fun2} ever returned, the \code{fun1}'s function body would be over
and code at the very end of the function (\emph{epilogue}) would set \code{SP}
to back to the value in \code{BP} like it was after the prologue.  Finally,
the function would restore the old (caller's) value of \code{BP} and invoke
\code{return} -- pop the address off the stack and jump there.

Conclusively, \code{BP} is somewhat static inside a function and changes only
by \codeasm{call}s and \code{return}s. Exactly, \code{BP} points to the
current frame plus 2 entries: return address and saved \code{BP}. This renders
\code{BP} to be only a convenient way for restoring \code{SP} in the epilogue.
Now, consider that that both \code{SP} and \code{BP} are usually stored in
registers. On \xesh, these are \codeasm{rsp} and \codeasm{rbp} (\codeasm{esp}
and \codeasm{ebp} on \xesl), respectively. If we could avoid use of \code{BP},
we would free up one register for arbitrary use. 

Convenient in the preceding paragraph means that if we properly track
throughout the function body what we perform with stack, we can easily restore
the stack state in the epilogue based on the operations. And pushes to and
pops from the stack are relatively easy to count, especially when the assembly
is generated by a compiler. So given the offset is known, by default, \gcc
with optimizations turned on do not use \code{BP} as a frame/base pointer
storage. Instead, it uses it as another registers. This behaviour can be
disabled by \code{-fno-omit-frame-pointer}, if needed, however.

\paragraph{Tail calls}

\subsection{Uses of Stack Traces} \label{sec:uses}
Crashes (incl.\ BUGs and WARNINGs) are obvious users of unwinders. They print
out the information so that developers can debug the crash. There are also
other users like \textsc{KMemLeak}, \textsc{KMemCheck}, or \textsc{LockDep}.
They all do use the stack traces to store information about locations where
memory accesses, allocations, and locks take place respectively.

We can consider the all above as debugging purposes. There is one more not so
obvious use of unwinders: \textit{Live patching}.

\section{Pre-rewrite Unwinding} \label{sec:pre}

\subsection{openSUSE \& SLE}

\section{Post-rewrite Unwinding} \label{sec:post}
In 4.9 by 7c7900f89770.

\subsection{openSUSE \& SLE}

\begin{figure}[tb]
  \begin{center}
    \begin{tabular}{|l||c|l|c|l|c|l|} \hline
      \multicolumn{1}{|c||}{\textbf{Union of Stack Traces}} &
	\multicolumn{6}{|c|}{\textbf{Unwinder}} \\ \cline{2-7}
      \multicolumn{1}{|c||}{\textbf{from All Unwinders}} &
	\multicolumn{2}{|c|}{\textbf{Guess}} &
	\multicolumn{2}{|c|}{\textbf{Frame Pointer}} &
	\multicolumn{2}{|c|}{\textbf{\orc}} \\ \hline\hline
      \code{crash\_now}		 & R & 0x94  & R & 0x99  & R & 0x94 \\
      \rowcolor[gray]{.9}
      \code{crash\_timer\_cb}	 &   &       &   & 0xe   &  & \\
      \code{call\_timer\_fn}	 & ? & 0x2e  &   & 0x33  &  & 0x2e \\
      \rowcolor[gray]{.9}
      \code{crash\_now}		 &   &       & ? & 0x170 &  & \\
      \code{run\_timer\_softirq} & ? & 0x40d &   & 0x414 &  & 0x40d \\
      \rowcolor[gray]{.9}
      \code{RANDOM}		 & ? & 0x1f  & ? & 0x2b  & ? & 0x1f \\
      \rowcolor[gray]{.9}
      \code{RANDOM}		 & ? & 0x38  & ? & 0x1d  & ? & 0x38 \\
      \rowcolor[gray]{.9}
      \code{RANDOM}		 & ? & 0x21  & ? & 0xc8  & ? & 0x21 \\
      \code{\_\_do\_softirq}	 & ? & 0xde  &   & 0xde  &  & 0xde \\
      \code{irq\_exit}		 & ? & 0xae  &   & 0xb6  &  & 0xae \\
      \code{smp\_apic\_timer\_interrupt} & ? & 0x39 &  & 0x3d &  & 0x39 \\
      \code{apic\_timer\_interrupt} & ? & 0x82 &  & 0x89 &  & 0x82
      \\ \hline\hline

      \multicolumn{7}{|c|}{--- IRQ ---} \\ \hline\hline

      \code{native\_safe\_halt}	 & ? & 0x2   & R & 0x6   & R & 0x2 \\
      \code{default\_idle}	 & ? & 0x1a  &   & 0x20  &   & 0x1a \\
      \rowcolor[gray]{.9}
      \code{arch\_cpu\_idle}	 &   &       &   & 0xf   &   & \\
      \rowcolor[gray]{.9}
      \code{default\_idle\_call} &   &       &   & 0x23  &   & \\
      \code{do\_idle}		 & ? & 0x16a &   & 0x178 &   & 0x16a \\
      \code{cpu\_startup\_entry} & ? & 0x5d  &   & 0x62  &   & 0x5d \\
      \rowcolor[gray]{.9}
      \code{rest\_init}		 &   &       &   & 0x84  &   & \\
      \code{start\_kernel}	 & ? & 0x422 &   & 0x429 &   & 0x422 \\
      \rowcolor[gray]{.9}
      \code{early\_idt\_handler\_array} & ? & 0x120 & ? & 0x120 & ? & 0x120 \\
      \rowcolor[gray]{.9}
      \code{x86\_64\_start\_reservations} &  &  &  & 0x29 &  & \\
      \code{x86\_64\_start\_kernel} & ? & 0x12c &  & 0x134 &  & 0x12c \\
      \code{secondary\_startup\_64} & ? & 0x9f &  & 0x9f &  & 0x9f \\ \hline
    \end{tabular}
  \end{center}
    \caption{Precision of unwinders demonstrated by \code{modprobe crasher
    call\_warn}. The table demonstrares the Guess, Frame Pointer and \orc
    unwinders. The question mark (?) means non-reliable. ``R'' means this
    value was in a register. The numerical columns are reported offsets
    against the start of the respective function. It can be also seen that
    the offsets of Guess and \orc unwinders are the same. I.e. \orc does not
    add any code anywhere, it relies only on out-of-band data, but still
    provides reliable stack traces.}
  \label{fig:instr}
\end{figure}

\section{Conclusion} \label{sec:concl}
Conc

\bibliographystyle{plain}
%\bibliography{paper}

\end{document}
